\def\udk{\ }



\def\tit{К 70-ЛЕТИЮ ЗАСЛУЖЕННОГО ДЕЯТЕЛЯ НАУКИ 

РОССИЙСКОЙ ФЕДЕРАЦИИ И.\,Н.~СИНИЦЫНА }



  \def\aut{С.\,Я.~Шоргин, Э.\,Р.~Корепанов, В.\,И.~Синицын}



\titel{\udk}{\tit}{\aut}



\def\razd{К 70-летию заслуженного деятеля науки 

РФ И.\,Н.~Синицына}



 \thispagestyle{empty}



\def\sect{С.\,Я.~Шоргин, Э.\,Р.~Корепанов, В.\,И.~Синицын}



\index{Шоргин С.\,Я.}

\index{Корепанов Э.\,Р.}

\index{Синицын В.\,И.}



%Институт проблем информатики РАН, Москва.



\begin{floatingfigure}{54mm} %fig5-1

\vspace*{6pt}

 \begin{center}

\mbox{%

\epsfxsize=50mm

\epsfbox{sok-1.eps}

}

\end{center}

\vspace*{9pt}

\end{floatingfigure}



    

       

\textbf{Игорь Николаевич Си-\linebreak ницын}, профессор, 

доктор технических наук~--- известный\linebreak ученый в 

области  прикладной механики и управления, 

при-\linebreak кладной ма\-те\-ма\-ти\-ки и информатики, 

основатель научной\linebreak шко\-лы в области 

стохастиче-\linebreak ских информационных техно-\linebreak логий и 

систем автоматизации, автор свыше 

500~научных трудов, 50~монографий, книг и\linebreak

30~изобретений. Имеет большой опыт работы в 

промышленности и высших учебных\linebreak заведениях.

      

      Игорь Николаевич Сини-\linebreak цын родился 

14~августа 1940~г.\ в Москве. Высшее 

образование получил в МВТУ им.\ Н.\,Э.~Баумана 

и МГУ им.\ М.\,В.~Ломоносова. 

Одновременно с учебой в МГУ начал работать в 

известном ракетно-космическом 

      на\-уч\-но-ис\-сле\-до\-ва\-тель\-ском институте, 

ныне Институте прикладной механики им.\ В.\,И.~Кузнецова (НИИПМ).

      

      В 1965~г.\ защитил кандидатскую диссертацию в МГУ, а в 1983~г.~--- докторскую 

в НИИПМ.

      

      В период с 1960 по 1983~гг.\ участвовал в разработке и испытаниях 

гироскопических командных приборов и информационно-из\-ме\-ри\-тель\-ных систем, вел 

на\-уч\-но-ис\-сле\-до\-ва\-тель\-скую  и инженерную деятельность, совмещая ее с педагогической: 

сначала в МВТУ им.\ Н.\,Э.~Баумана, затем в Военно-воздушной 

инженерной академии им.\ профессора Н.\,Е.~Жуковского (ВВИА). В~дальнейшем 

работал на факультете авиационного вооружения\linebreak ВВИА, где занимался подготовкой 

авиационных инженеров и космонавтов, принимал участие в разработке и испытаниях 

специальной техники.

      

      В 1984~г.\ И.\,Н.~Синицын был переведен во вновь организованный Институт 

проблем информатики АН СССР (ныне ИПИ РАН) для проведения работ в области 

специальных системных применений ЭВМ новых поколений.

      

      В настоящее время И.\,Н.~Синицын работает в ИПИ РАН заведующим отделом 

<<Стохастические проблемы информатики и управления>>. С~1987~г.\ 

И.\,Н.~Синицын~--- профессор МАИ.\linebreak В~разные годы И.\,Н.~Синицын был 

заместителем Генерального конструктора, заместителем главного конструктора и главным 

конструктором ряда автоматизированных и информационных\linebreak систем специального 

назначения.

     

     Его основные научные труды относятся к следующим об\-ластям:

     \begin{itemize}

     \item статистическая теория информационных технологий и автоматизированных 

систем;

     \item прецизионные информационно-измерительные технологии и системы для 

научных исследований и специального назначения;

     \item информационно-аналитические технологии и системы поддержки принятия 

решений для информатизации высших органов государственной власти РФ, федеральных 

ве-\linebreak домств и~др.

     \end{itemize}

     

     В области статистической теории информационных тех\-но\-логий и 

автоматизированных систем И.\,Н.~Синицыну принад\-лежат фундаментальные 

результаты, относящиеся к теории канонических представлений случайных функций в 

сложных\linebreak стохастических системах (СтС), в том числе в СтС с распределенными 

параметрами и случайной структурой (1964--1983~гг.). И.\,Н.~Синицын распространил 

методы теории СтС на системы, описываемые дифференциальными уравнениями со 

случайными функциями состояния, уравнениями в гильбертовых и банаховых 

пространствах. Им разработаны эффективные вычислительные методы нахождения 

распределений, основанные на параметризации и позволяющие радикально сократить 

количество уравнений для параметров распределений, а также новые вы\-чис\-ли\-тель\-ные 

методы статистического анализа и синтеза, допускающие эффективное оценивание 

точности и ориентированные на параллельные статистические вычисления. 

     

И.\,Н.~Синицын разработал методы нахождения точных выражений для 

распределений с инвариантной мерой, обнаружил ряд новых классов точных 

распределений (1979--1993~гг.). Им получены фундаментальные результаты в области 

нелинейной условно оптимальной и субоптимальной фильтрации в реальном масштабе 

времени. Важные результаты получены И.\,Н.~Синицыным по теоретико-групповым 

методам анализа и синтеза\linebreak  автоматизированных систем (1983--2000~гг.). Им разработана 

статистическая теория катастрофоустойчивости автоматизированных систем высокой 

точности и доступности. 

     

И.\,Н.~Синицын~--- основоположник стохастических ин\-фор-\linebreak мационных технологий 

оперативной обработки информации,\linebreak контроля и мониторинга автоматизированных 

систем, а также технологии стохастического управления информационными активами. 

Под его руководством и при его непосредственном учас\-тии созданы инструментальные 

средства для статистического анализа, моделирования и синтеза систем: 

     проб\-лем\-но-ориен\-ти\-ро\-ван\-ные диалоговые системы и библиотеки <<СтС-Анализ>>\linebreak 

(1989--1990, 2004--2007), <<СтС-Фильтр>> (1991, 2004--2007),\linebreak <<СтС-Модель>> (1993, 

2007), Nalib (1991--1993), <<TransStatLib>> (1991--1993),  ASAP-ALPHA-97, 

<<Безопасность и надежность>>\linebreak (1991--1995~гг.), <<Здоровье РФ>> (1992--1996, 1999, 

2003) и~др. В~последние годы им разработаны эффективные символьные методы анализа 

и синтеза СтС. В~2005--2007~гг.\ создано и внедрено специализированное программное 

обеспечение <<СтС-СМА>> и\linebreak <<СтС-ИТКР>> (2007--2010). 

     

     Его книги~--- <<Стохастические дифференциальные системы. Анализ и 

фильтрация>> (1985, 1987, 1990), <<Лекции по функциональному анализу и его 

приложениям>> (1999), <<Теория стохастических систем>> (2000, 2001, 2004) (совместно 

с В.\,С.~Пугачевым), а также <<Фильтры Калмана и Пугачева>> (2006, 2007) и 

<<Канонические представления случайных функций и их применение в задачах 

компьютерной поддержки научных исследований>> (2009) широко известны в России и 

за рубежом.

     

     В области прецизионных информационно-измерительных\linebreak технологий и 

автоматизированных систем для научных исследований и специального назначения 

И.\,Н.~Синицыным впервые разработана теория ряда информационно-измерительных 

систем в условиях случайных динамических возмущений, открыт ряд новых 

статистических динамических эффектов (выбросы разных типов, флуктуационные уходы, 

накопление возмущений и~др.). Ему принадлежат первые работы по статистической 

динами-\linebreak ке ко\-мандно-из\-ме\-ри\-тель\-ных гироскопических приборов, акселе-\linebreak рометров, 

градиентометров и метрологических систем высочай\-шей точности, информационной 

теории и методам измерений,\linebreak калибровок, ускоренных испытаний в экстремальных 

условиях, а также статистического и полунатурного моделирования. Под его 

руководством и при его непосредственном участии разработано и внедрено несколько 

поколений серийных систем (в~том чис\-ле со встроенными устройствами для снижения 

возмуще-\linebreak ний), обладающих уникальными характеристиками (1959--\linebreak 1983~гг.). 

И.\,Н.~Синицын принимал непосредственное участие в определении технической 

политики в области новой специальной техники (1979--1983~гг.). 

     

     С именем И.\,Н.~Синицына связано создание концепций автоматизации научных 

исследований в РФ, в первую очередь основанных на средствах массовой 

вычислительной техники (1983--1988~гг.). 

     

     Под руководством и при участии И.\,Н.~Синицына сформу-\linebreak лированы принципы 

создания микровидеосистем, создан ряд\linebreak базовых систем. Они внедрены в МВД и 

Минздраве РФ (1983--1988~гг.). Полученные результаты получили развитие в 

автоматизированных системах метрологического обеспечения,\linebreak видео\-контроля и 

биометрических системах. Под руководством И.\,Н.~Синицына разработаны принципы 

построения и архитектуры вычислительных систем командных пунктов, а также новые 

быст\-рые методы и алгоритмы обработки изображений, обла\-да\-ющих сильной 

пространственно-временной де\-фор\-ма\-ци\-ей.\linebreak В~интересах автоматизации астрометрических 

научных исследований были созданы комплекс моделей, алгоритмов и специального 

программного обеспечения и информационные ресурсы для нестандартной 

интегрированной обработки параметров вращения Земли по фундаментальной проблеме 

<<Статистическая динамика вращения Земли>>. В~2005--2010~гг.\ им впервые 

обнаружен ряд новых эффектов (автоколебания полюса Земли на чандлеровской частоте, 

параметрическая стабилизация чандлеровских колебаний, нелинейные флуктуационные 

дрейфы не\-ста\-биль\-ности вращения Земли и~др.).

     

     В области информационно-аналитических технологий и автоматизированных систем 

поддержки принятия решений для\linebreak 

информатизации высших органов государственной 

власти, федеральных ведомств и~др.\ И.\,Н.~Синицыным разработан и внедрен\linebreak ряд 

базовых информационных технологий (обработка информации от независимых 

источников, формирование и хранение больших баз электронных образов, управление 

информационными активами и~др.). Сформулированы принципы и разработаны базовые 

системотехнические решения для ряда крупномасштабных автоматизированных 

информационных и ин\-фор\-ма\-ци\-он\-но-управ\-ля\-ющих систем специального назначения, 

высокой точ-\linebreak ности и доступности. 

     

     Синицын~И.\,Н. отдает много сил подготовке научно-пе\-да\-го\-ги\-че\-ских кадров. Под 

его руководством защищено свыше 25~кандидатских и докторских диссертаций. Он 

является заместителем Председателя и руководителем специальной секции\linebreak Ученого 

совета ИПИ РАН, состоит членом ряда специализированных диссертационных советов, 

членом экспертного совета\linebreak РФФИ, заместителем главного редактора журналов 

<<Наукоемкие технологии>> и <<Системы высокой доступности>>, а также членом 

редколлегий ряда отечественных и международных журналов. И.\,Н.~Синицын 

возглавляет Комиссию Минобразования РФ по информатике в военных учебных 

заведениях.

     

     И.\,Н.~Синицын пользуется широким международным авторитетом. Он принимал 

активное участие в организации международного сотрудничества стран СЭВ в области 

ЭВМ новых поколений. Ряд его книг издан за рубежом. В 1990--1994~гг.\ он был 

директором Российско-Французского центра <<Эвклид>>; в период с 1985 по 2006~гг.\ 

был экспертом фонда INTAS; с 1983~г.\ по настоящее время является членом 

программных комитетов ряда международных форумов в области теоретической и 

прикладной информатики.

     

     \bigskip

     Коллектив Института проблем информатики РАН сердечно поздравляет Игоря 

Николаевича Синицына с 70-летием! Мы желаем юбиляру крепкого здоровья, новых 

научных достижений, счастья и успехов.

